% Generated draft: requires mathematical and source review.
\documentclass[11pt]{article}
\usepackage[margin=25mm]{geometry}
\usepackage{fontspec}
\setmainfont{Latin Modern Roman}
\usepackage{amsmath,amssymb,mathtools,unicode-math}
\setmathfont{Latin Modern Math}
\usepackage{enumitem,needspace,xurl,hyperref}
\setlength{\parindent}{0pt}
\setlength{\parskip}{5pt}
\setlength{\emergencystretch}{3em}
\newcommand{\N}{\mathbb{N}}
\newcommand{\Z}{\mathbb{Z}}
\newcommand{\Q}{\mathbb{Q}}
\newcommand{\R}{\mathbb{R}}
\newcommand{\C}{\mathbb{C}}
\newcommand{\F}{\mathbb{F}}
\newcommand{\A}{\mathbb{A}}
\newcommand{\PP}{\mathbb{P}}
\newcommand{\E}{\mathbb{E}}
\newcommand{\eps}{\varepsilon}
\newcommand{\dd}{\,\mathrm{d}}
\begin{document}
\section*{UnsolvedMath ID: 30001745}
\textbf{Source problem code:} catalogue entry
\subsection{Definitions and mathematical statement}
A polynomial vector field of degree at most two on an affine space over a field \(K\) of characteristic zero can be written as
\[
\dot x=f(x)=Q(x)+Bx+c,\qquad x\in K^n,
\]
where \(Q\colon K^n\to K^n\) is homogeneous quadratic, \(B\) is linear, and \(c\in K^n\). Its symmetric bilinear polarization is
\[
q(u,v)=\frac{Q(u+v)-Q(u)-Q(v)}{2},\qquad q(x,x)=Q(x).
\]
In a standard step-size normalization, the Hirota--Kimura discretization is the rational map \(\Phi_h\colon x\dashrightarrow x^+\) implicitly defined by
\[
\frac{x^+-x}{h}
=q(x,x^+)+\frac12 B(x+x^+)+c,
\qquad h\ne0.
\]
This equation is linear in \(x^+\); it defines a rational map on the locus where the corresponding coefficient matrix is invertible. Interchanging \(x\) and \(x^+\) replaces \(h\) by \(-h\), giving the rational inverse wherever defined. This is a conventional description of the named discretization; the supplied record gives no formula or step-size normalization and does not specify whether ``quadratic'' means homogeneous quadratic or degree at most two.

For orientation, in a usual complex-algebraic interpretation, algebraic complete integrability of a continuous system involves algebraic first integrals whose generic common level varieties are Zariski-open subsets of abelian varieties, with the vector field extending to a translation-invariant vector field on those abelian varieties. An abelian variety is a connected projective algebraic group. A corresponding interpretation for a discrete system requires invariant level varieties with such compactifications and dynamics extending to translations on them. These descriptions indicate the geometric content of the question; they do not supply all hypotheses of a unique definition. The source record does not specify the phase space, the number and independence conditions of the first integrals, any Hamiltonian or Poisson structure, or the precise discrete integrability criterion.

\textbf{Question supplied by the source.} Does every algebraically completely integrable system with a quadratic vector field remain algebraically completely integrable after applying its Hirota--Kimura discretization? Thus the proposed universal implication is
\[
\begin{gathered}
f\text{ is a quadratic vector field and its continuous system}\
\text{is algebraically completely integrable}
\end{gathered}
\quad\Longrightarrow\quad
\Phi_h\text{ is algebraically completely integrable}.
\]
The quantifier ranges over every system in the intended class. A fully specified assertion still requires a choice of that class and its integrability definitions. In particular, the supplied record does not determine whether the conclusion is required for generic nonzero \(h\), for every nonzero \(h\) for which a birational map exists, or in another step-size regime. Exceptional points where the rational map is undefined must also be distinguished from its generic invariant varieties. These unresolved qualifications are retained rather than assigned additional hypotheses.
\subsection{Short English statement}
The question asks whether Hirota--Kimura discretization preserves algebraic complete integrability for every quadratic vector field whose continuous system has that property. The source leaves the precise phase space, integrability conventions, and step-size quantifiers unspecified. The displayed discretization uses a standard convention. This is a statement-only draft requiring human mathematical review.
\subsection{Sources}
\begin{itemize}
\item Local UnsolvedMath registry: supplied complete record for permanent ID \(30001745\), titled ``Integrability Under Hirota--Kimura Discretization,'' in the Oberwolfach Reports Open Problems collection. Exact source problem code: \texttt{catalogue entry}. The supplied proposed year is 2011; no individual proposer is named.
\item Supplied report citation: \emph{Geometric Numerical Integration} (2011), Oberwolfach Reports, pp.\ 825--900, DOI \texttt{10.4171/owr/2011/16}. Supplied source URL and literature-source link labelled ``original OWR report'': \url{https://doi.org/10.4171/owr/2011/16}.
\item Supplied literature-source link labelled ``Hirota--Kimura integrability survey'': \url{https://arxiv.org/abs/1008.1040}. No authors, publication title, or specific supporting passage are supplied for this link.
\item Literature checked at: \texttt{2026-08-21}. This is supplied dated curation, not an independently verified assessment. It labels the problem open and states that the general conjecture remains unresolved, that surveys report indications against it, and that no rigorous counterexample settling the universal claim is reported. The background expressly requests independent verification of time-sensitive status and cited claims.
\item The registry labels the statement \texttt{corrected\_verified} and supplies the verification note that the official report confirms the integrability question. These are supplied source-verification claims; no independent source inspection is asserted here.
\end{itemize}
\end{document}
